\documentclass[11pt,a4paper]{article}
\usepackage[T1]{fontenc}
\usepackage[utf8]{inputenc}
\usepackage{lmodern}
\usepackage{amsmath,amssymb}
\usepackage[margin=25mm]{geometry}
\usepackage{longtable,booktabs,array,calc}
\usepackage[hidelinks]{hyperref}
\hypersetup{pdftitle={Operational spin states, reversible interaction and finite closure},pdfauthor={navstokgap research working notes}}
\newcommand{\tightlist}{\setlength{\itemsep}{0pt}\setlength{\parskip}{0pt}}
\setlength{\emergencystretch}{3em}
\setcounter{secnumdepth}{0}
\begin{document}
\hypertarget{operational-spin-states-reversible-interaction-and-finite-closure}{%
\section{Operational spin states, reversible interaction and finite
closure}\label{operational-spin-states-reversible-interaction-and-finite-closure}}

For classical unit orientations with all Borel preparations and
restricted affine product tests, local rotations close on a
finite-dimensional ball state space. The complete reversible interaction
\(H=J u_zv_z\), however, does not descend to its moment tensor, and no
finite enlargement by bounded observables repairs exact descent on any
nonzero time interval. These are explicit classical countermodels and
proofs. They identify the missing operational closure premise in
source-backed reconstructions of quantum composition; they do not select
a positive universal action parameter.

The model below keeps its preparations, effects and transformations
explicit. It incorporates C125 and C127--C129 and the associated
source-premise map. Historical coordinator/model checks are internal.
The reconstruction theorems are cited with their separate hypotheses and
reading levels; their complete proofs have not been independently
audited here.

\hypertarget{a-ball-of-local-states-and-a-minimal-composite}{%
\subsection{A ball of local states and a minimal
composite}\label{a-ball-of-local-states-and-a-minimal-composite}}

Let the hidden state be \(u\in S^2\), with every Borel probability
measure \(\mu\) admitted as a preparation. A local measurement is a
finite family

\[e_i(u)=a_i+b_i\cdot u,\qquad a_i\ge|b_i|,\qquad
\sum_i a_i=1,\quad\sum_i b_i=0.\]

Its probabilities depend exactly on \(r=\int u\,d\mu\). Every point of
the closed unit ball is achieved by mixing opposite orientations, and
the tests \((1\pm n\cdot u)/2\) determine \(r\). Thus operational
equivalence is equality of \(r\), pure states are the unit sphere, and
the unnormalized state cone has linear dimension four. If a measurement
distinguishes \(k\) states perfectly, then each successful outcome obeys
\(1=a_i+b_i\cdot r_i\le2a_i\), hence \(k\le2\). Opposite poles attain
capacity two. Local \(SO(3)\) rotations descend exactly to
\(r\mapsto Rr\) and connect all pure states continuously.

For \(n\) components, admit all joint Borel measures on \((S^2)^n\),
product local tests, mixtures and coarse-grainings, local rotations and
permutations. Put \(v(u)=(1,u)\). All test probabilities depend on

\[z=\int v(u_1)\otimes\cdots\otimes v(u_n)\,d\mu.\]

Product tests span these coordinates. The composite is exactly the
convex hull of product ball states, with linear dimension \(4^n\). Each
product effect is at most \(2^n\) times its value at the product of zero
vectors. This survives mixing and coarse-graining; a test distinguishing
\(k\) states must spend at least \(2^{-n}\) of its normalization per
successful outcome, so \(k\le2^n\). Product pole tests attain \(2^n\).

These are stipulated terminal measurements, not unrestricted
hidden-state readouts or invasive instruments. If every hidden
orientation event were measurable, the operational state would
distinguish the full probability measures. The model is valid with its
local transformations; additional interactions must be tested for
descent.

\hypertarget{what-reconstruction-premises-add}{%
\subsection{What reconstruction premises
add}\label{what-reconstruction-premises-add}}

This ball model satisfies local pure-state connectivity but not the
whole axiom package of either reconstruction below:

\begin{longtable}[]{@{}
  >{\raggedright\arraybackslash}p{(\columnwidth - 4\tabcolsep) * \real{0.3333}}
  >{\raggedright\arraybackslash}p{(\columnwidth - 4\tabcolsep) * \real{0.3333}}
  >{\raggedright\arraybackslash}p{(\columnwidth - 4\tabcolsep) * \real{0.3333}}@{}}
\toprule\noalign{}
\begin{minipage}[b]{\linewidth}\raggedright
Source and selected reading
\end{minipage} & \begin{minipage}[b]{\linewidth}\raggedright
Additional premises
\end{minipage} & \begin{minipage}[b]{\linewidth}\raggedright
Classical comparison
\end{minipage} \\
\midrule\noalign{}
\endhead
\bottomrule\noalign{}
\endlastfoot
Hardy, version 4, §§1, 4, 7; \href{../references/batches/B67.md}{B67},
passages & Probability, simplicity, subspace and composition axioms,
plus continuous reversible pure-state connectivity & A finite simplex
has isolated vertices and only vertex permutations; the orientation ball
evades this particular failure \\
Chiribella--D'Ariano--Perinotti, version 3, §II and §§III.1--III.2; B67,
passages & Causality, perfect distinguishability, ideal compression,
local distinguishability, pure conditioning and purification with
essential uniqueness & Pure states of the minimal classical composite
have pure marginals, so a mixed local state has no purification \\
\end{longtable}

The purification assertion follows directly. The product-pure set is
compact; every extreme point of its convex hull is product-pure by a
finite convex representation. Conversely a unit-vector marginal forces
that orientation almost surely, successively in every factor, so
product-pure points are extreme. Every marginal of a pure composite is
therefore pure. The state \(r=0\) has no pure extension by any number of
these components.

Hardy's subspace condition also fails explicitly. In two components,
coarse-grain the local pole tests along \(n\) to equal signs:

\[e_{\rm same}(u,v)=\tfrac12[1+(n\cdot u)(n\cdot v)].\]

Probability one requires \((n\cdot u)(n\cdot v)=1\) almost surely,
leaving only \((n,n)\) and \((-n,-n)\). The corresponding capacity-two
face is a segment, whereas an elementary capacity-two system here is a
three-dimensional ball. Local continuity and multiplicative
capacity/dimension therefore do not select quantum composition.

Hardy's moving-ball comparison motivates closure throughout a
transformation: coarse-graining intermediate positions into a bit does
not make their reduced evolution reversible. Likewise a definite
mechanical microstate is not by itself a CDP operational purification.
The full reconstruction packages remain conditional inputs. Originals:
\href{https://arxiv.org/html/quant-ph/0101012v4}{Hardy} and
\href{https://arxiv.org/html/1011.6451v3}{CDP}. Only the indicated
passages and hypothesis comparisons were read; no claim is made about
all continuous classical phase spaces or the infinite-dimensional
extensions.

\hypertarget{a-complete-classical-interaction-fails-moment-descent}{%
\subsection{A complete classical interaction fails moment
descent}\label{a-complete-classical-interaction-fails-moment-descent}}

Fix spin magnitudes \(S_A,S_B>0\) with action units and \(J>0\) with
energy units. On \(S^2\times S^2\) use

\[\{u_i,u_j\}=S_A^{-1}\epsilon_{ijk}u_k,\qquad
\{v_i,v_j\}=S_B^{-1}\epsilon_{ijk}v_k,\qquad \{u_i,v_j\}=0,
\qquad H=J u_zv_z.\]

With \(\dot f=\{f,H\}\) and \(\Omega_A=J/S_A\), \(\Omega_B=J/S_B\),

\[\dot u_x=-\Omega_Au_yv_z,\quad\dot u_y=\Omega_Au_xv_z,\quad\dot u_z=0,\]
\[\dot v_x=-\Omega_Bv_yu_z,\quad\dot v_y=\Omega_Bv_xu_z,\quad\dot v_z=0.\]

Norms and energy are conserved. Smoothness on compact phase space makes
the flow complete and reversible, with inverse at negative time. Its
frequencies depend on the other spin, so this is an actual hidden
interaction.

Prepare the exact pair

\[\mu_A=\delta_{e_y}(du)\tfrac12(\delta_{e_z}+\delta_{-e_z})(dv),\qquad
\mu_B=\delta_{e_y}(du)\tfrac12(\delta_{e_x}+\delta_{-e_x})(dv).\]

Both have \(\mathbb E[u]=e_y\), \(\mathbb E[v]=0\),
\(\mathbb E[uv^T]=0\), and microscopic energy identically zero. They
agree on all original terminal tests. Since \(u_z=0\), \(v\) stays
constant. In the first preparation,

\[u(t)=(-\operatorname{sign}(v_z)\sin\Omega_At,\cos\Omega_At,0),\qquad
\mathbb E_{\mu_A}[u_x(t)v_z(t)]=-\sin\Omega_At.\]

The second retains \(u=e_y\) and has zero cross expectation. The allowed
product effect \((1+u_x)(1+v_z)/4\) consequently gives

\[p_A(t)=\tfrac14(1-\sin\Omega_At),\qquad p_B(t)=\tfrac14.\]

These differ at every sufficiently small positive time. No single-valued
update, linear or nonlinear, on the initial moment tensor can represent
this flow for all preparations. Both evolved measures still give valid
minimal-composite states: the failure is descent to the quotient. Its
missing moment is exposed by

\[\frac d{dt}\mathbb E[u_xv_z]=-\Omega_A\mathbb E[u_yv_z^2].\]

Source/proof context: \href{../references/batches/B72.md}{B72}, written
spin conventions and internal comparison. The supplied \(S_A,S_B\) are
not derived Planck constants.

\hypertarget{no-finite-observable-enlargement-repairs-this-interaction}{%
\subsection{No finite observable enlargement repairs this
interaction}\label{no-finite-observable-enlargement-repairs-this-interaction}}

Let \(V\) be a finite-dimensional real space of bounded Borel functions
on \(X=S^2\times S^2\), compared pointwise, containing \(1\) and the
original terminal effects. Exact closure would require
\(U_tV\subset V\), where \(U_tf=f\circ\Phi_t\), for every \(t\) in a
nonzero interval. A forward interval is sufficient for the
contradiction. The excluded case \(J=0\) has static finite closure.

In the complexification, \(w=u_x+iu_y\in V_{\mathbb C}\) obeys exactly

\[U_tw=w\exp(i\Omega_Atv_z).\]

For distinct \(t_1,\ldots,t_m\), a relation \(\sum_j c_jU_{t_j}w=0\)
restricted to \(u=e_x\) and \(v=(\sqrt{1-z^2},0,z)\), \(-1<z<1\), gives

\[\sum_jc_j\exp(i\Omega_At_jz)=0.\]

Derivatives at \(z=0\) of orders \(0,\ldots,m-1\) form a Vandermonde
system with determinant \(\prod_{j<l}(i\Omega_At_l-i\Omega_At_j)\ne0\).
Every coefficient vanishes. Arbitrarily many time translates are
therefore independent, contradicting finite dimension. No regularity of
added observables is needed. The smooth infinitesimal hierarchy is
equivalently \(L^nw=(i\Omega_A)^nwv_z^n\), but the exact time-orbit
proof establishes closure failure without a truncation argument.

Nor can a nonlinear update of finitely many expectations evade it. If
the expectation of \(g\) is determined by those of
\(f_0=1,f_1,\ldots,f_d\) for \textbf{every} probability measure, then
\(g\) lies in their real span. Otherwise finite evaluation linear
algebra supplies points \(x_j\) and weights \(a_j\) with
\(\sum_ja_jf_k(x_j)=0\) for every \(k\), but \(\sum_ja_jg(x_j)\ne0\). If
no such relation existed, \(g(x)\) would define a linear functional on
the span of the evaluation vectors and hence be a combination of the
\(f_k\). Since \(\sum_ja_j=0\), positive and negative weights have equal
nonzero mass; normalizing them gives two admitted atomic preparations
with equal retained expectations and unequal \(g\) expectation. Apply
this to \(g=U_tf_k\).

Thus exact finite expectation coordinates require an invariant finite
observable span, already excluded above. Arbitrary encodings of whole
measures without affine probabilistic mixing are outside this theorem.
Infinite observable dimension does not by itself imply infinite
perfect-distinguishability capacity or quantum structure. Restricted
preparations and finite accuracy change its hypotheses. Source/proof
context: \href{../references/batches/B74.md}{B74}, written proof and
literature boundary.

\hypertarget{reversible-operational-interaction-the-source-theorem}{%
\subsection{Reversible operational interaction: the source
theorem}\label{reversible-operational-interaction-the-source-theorem}}

A different premise demands a connected reversible \textbf{linear} group
on the fixed operational tensor coordinates, containing local rotations
and a genuine interaction. Under \(\rho(r)=(I+r\cdot\sigma)/2\),
\(E=aI+b\cdot\sigma\), local probabilities are
\(\operatorname{tr}(E\rho)=a+b\cdot r\). The tensor coordinates become
Hermitian matrices and the minimal composite becomes exactly separable
density matrices. No entangled state was inserted by this coordinate
identification.

Theorem 1 of de la Torre--Masanes--Short--Müller assumes valid
probabilities for every product preparation, group element and product
test, with the group containing local unitaries. A larger connected
group contains, on a pair with fixed product spectators, an entangling
unitary action or its partial-transpose conjugate. The first leaves the
minimal composite by creating entanglement. For the second, choose a
product input whose partially transposed input is entangled by that
unitary: the output is the partial transpose of a pure entangled state.
Its Schmidt eigenvalues include \(\pm s_1s_2\), \(s_1s_2>0\), so it
cannot be separable positive semidefinite. Either branch excludes an
interacting extension preserving the minimal cone.

Theorem 2 additionally uses interchangeable identical systems and
three-system consistency to eliminate the partial-transpose alternative.
Entangling-gate universality and ancilla preparation, joint evolution,
measurement and discard closure supply the full quantum preparation,
measurement and completely positive channel conclusion. Those
assumptions must accompany that conclusion. The theorem concerns finite
numbers of identical qubits, not arbitrary local state spaces.

The source statements, assumptions and these consequences were read at
passage level in \href{../references/batches/B70.md}{B70}. The technical
Lie-algebra and gate-universality proofs are borrowed results, not
independently audited proofs of this repository. The exact hidden
Hamiltonian above does not contradict the theorem: it lacks reversible
operational descent, which no finite moment repair supplies under its
preparation hypotheses.

\hypertarget{the-checked-first-order-generator-restriction}{%
\subsection{The checked first-order generator
restriction}\label{the-checked-first-order-generator-restriction}}

A specific source step does have a full written verification here
(C129). Let \(X\) be a real \(4^n\times4^n\) matrix with \(e^{tX}\)
admissibly reversible for a two-sided neighbourhood of zero.
Independently available unit Bloch vectors in each input and effect give

\[p_{b,a}(t)=2^{-n}v(b)^Te^{tX}v(a)\in[0,1].\]

The expansion is

\[p(t)=p(0)+2^{-n}t\,v(b)^TXv(a)
 +2^{-n}\tfrac{t^2}2v(b)^TX^2v(a)+O(t^3).\]

An orthogonal pair \(b_1=-a_1\) gives \(p(0)=0\), so the two-sided
minimum requires \(v(-a_1,b_2,\ldots)^TXv(a)=0\) and the corresponding
\(X^2\) quadratic form nonnegative. With \(b_k=a_k\) at every site,
\(p(0)=1\) and \(v(a)^TX^2v(a)\le0\). These necessary derivative signs
are not sufficient finite-time positivity conditions. A forward
semigroup alone permits a nonzero first derivative at a boundary.

Fix spectator input and output indices. The vectors
\(v(e_i)\pm v(-e_i)\) span \(\mathbb R^4\), so independent product
spanning implies the local \(4\times4\) slice \(M\) satisfies
\(v(-a)^TMv(a)=0\) for every \(|a|=1\). Writing
\(M=\bigl(\begin{smallmatrix}c&r^T\\s&D\end{smallmatrix}\bigr)\), this
says \(c+(r-s)\cdot a-a^TDa=0\). Values at \(a,-a\) give \(r=s\); at
\(e_i\) they give \(D_{ii}=c\); at \((e_i+e_j)/\sqrt2\) they give
\(D_{ij}+D_{ji}=0\). Thus exactly

\[M=cI_4+\begin{pmatrix}0&r^T\\r&K\end{pmatrix},\qquad K^T=-K.\]

These slices form the seven-dimensional space
\(\mathcal L=\mathcal A\oplus\mathcal B\oplus\mathcal I\), with three
antisymmetric spatial directions, three symmetric time--space directions
and one scalar. At each site \(X\) lies in the tensor space with that
site restricted to \(\mathcal L\). Expand in a basis adapted to
\(\mathcal L\) and a complement; intersecting all site restrictions
leaves

\[X\in\mathcal L^{\otimes n}.\]

This is a necessary ambient linear space, not the Lie algebra of an
admissible group. Normalization, the second-order inequalities and later
source steps remain. Neither local rotation averaging nor
ancilla/discard closure was needed for this first-order restriction.
Source: the
\href{../docs/batches/B70/de-la-Torre-Masanes-Short-Mueller-2011.pdf}{held
original}, PDF pp.~3--4, equations (4), (6), (8)--(14);
\href{../reviews/reversible-generator-B77.md}{B77 internal check}.

\hypertarget{consequence-for-state}{%
\subsection{Consequence for STATE}\label{consequence-for-state}}

Finite operational capacity and smooth classical reversibility coexist
with failure of composite descent. Exact closure for the specified
interacting spin model requires infinite observable dimension or a
changed preparation, effect or dynamics class. Quantum reconstruction
needs its stated additional composition premises. All probability
coordinates above are dimensionless; matching their transformations to
mechanical time, energy and phase remains separate from determining
quantum structure or a universal physical \(\hbar\).
\end{document}
